\documentclass[11pt,a4paper]{article}

\usepackage[margin=1in]{geometry}
\usepackage{amsmath,amssymb,amsthm}
\usepackage{mathtools}
\usepackage{enumerate}
\usepackage{hyperref}
\usepackage{cleveref}
\usepackage[ruled,vlined,linesnumbered]{algorithm2e}
\usepackage{tikz}
\usetikzlibrary{arrows.meta,positioning}

\newtheorem{theorem}{Theorem}[section]
\newtheorem{lemma}[theorem]{Lemma}
\newtheorem{proposition}[theorem]{Proposition}
\newtheorem{corollary}[theorem]{Corollary}
\newtheorem{definition}[theorem]{Definition}
\newtheorem{remark}[theorem]{Remark}
\newtheorem{example}[theorem]{Example}

\DeclareMathOperator{\proj}{proj}
\DeclareMathOperator{\AM}{AM}
\DeclareMathOperator{\HM}{HM}
\DeclareMathOperator{\diag}{diag}
\DeclareMathOperator{\tr}{tr}
\DeclareMathOperator*{\argmax}{arg\,max}
\DeclareMathOperator*{\argmin}{arg\,min}
\newcommand{\R}{\mathbb{R}}
\newcommand{\E}{\mathbb{E}}
\newcommand{\Prob}{\mathbb{P}}
\newcommand{\cF}{\mathcal{F}}
\newcommand{\cG}{\mathcal{G}}
\newcommand{\cS}{\mathcal{S}}
\newcommand{\cA}{\mathcal{A}}
\newcommand{\cO}{\mathcal{O}}
\newcommand{\cC}{\mathcal{C}}
\newcommand{\cJ}{\mathcal{J}}
\newcommand{\ip}[2]{\langle #1, #2 \rangle}
\newcommand{\norm}[1]{\|#1\|}
\newcommand{\Lself}{L_{\text{self}}}
\newcommand{\Lch}{L_{\text{ch}}}
\newcommand{\Lcomm}{L_{\text{comm}}}

\title{Evidence-Weighted Multi-Agent Debate\\
for Scalable Oversight:\\
An Application of the $\Omega$-Framework}
\author{Eugene Shcherbinin\\
\small BSc Mathematics, London School of Economics}
\date{March 2026}

\begin{document}
\maketitle

\begin{abstract}
We apply the $\Omega$-framework (evidence-weighted multi-agent policy gradient)
to the problem of AI safety via debate. Two debaters---a proponent and an
opponent---compete over a truth claim while cooperating with a
bounded-rationality judge. We formalize debate as a multi-agent stochastic
game and derive three main results. First, debate accuracy is bounded by
$\min(C, \Lself^P, \Lself^O)$---the minimum of the judge's channel capacity
and the debaters' self-knowledge---so that models which ``vibe'' (perform
well through implicit pattern-matching without self-understanding) cannot be
effectively overseen through debate regardless of debate length. Second,
evidence-weighted debate, where the judge applies Keynesian weights to
arguments based on debater track records, achieves an $\HM(V)/\AM(V)$
improvement in accuracy over uniform-weight debate. Third, when arguments
are sparse ($k \ll d$ relevant argument types among $d$ possible), the judge
needs only $O(k \log d)$ rounds to identify the correct claim, connecting
concise reasoning to the blessing of dimensionality. We extend these
results to $N$-agent coalition debate and discuss implications for scalable
oversight: debate is a reliable oversight mechanism only for models with
sufficient self-knowledge, providing a quantitative criterion for when
debate-based alignment succeeds or fails.
\end{abstract}

\tableofcontents

%% ========================================================================
\section{Introduction}
\label{sec:intro}
%% ========================================================================

AI safety via debate \cite{irving2018} proposes that two AI systems arguing
before a human judge can produce aligned outcomes even when the judge cannot
independently verify the debaters' claims. The intuition is game-theoretic:
in a zero-sum debate, the optimal strategy is truth-telling, because a
truthful debater can always refute a lying opponent. Empirical work
\cite{khan2024,bowman2022} has partially validated this intuition but
also revealed limitations: debate quality depends on the debaters' ability
to \emph{explain} their reasoning, not just to be correct.

This paper connects debate to the $\Omega$-framework for multi-agent
learning \cite{shcherbinin2026ewpg,shcherbinin2026coop,shcherbinin2026bless},
which formalizes five components: evidence-weighted gradients, opponent
shaping, cooperative communication with lossy channels, the self-knowledge
bottleneck, and the blessing of dimensionality for sparse policies. Debate
maps onto this framework with striking precision:

\begin{center}
\begin{tabular}{lll}
\textbf{$\Omega$-Framework} & & \textbf{Debate Setting} \\
\hline
Policy $\pi_i$ & $\mapsto$ & Argument strategy \\
Observation $o_t$ & $\mapsto$ & Previous arguments \\
Evidence weight $w_i = V_{\min}/V_i$ & $\mapsto$ & Debater track record \\
Communication channel & $\mapsto$ & Natural language (lossy) \\
Self-knowledge loss $\Lself^{(i)}$ & $\mapsto$ & Ability to explain reasoning \\
$k$-sparse policy & $\mapsto$ & Concise argument (few key points) \\
Coalition & $\mapsto$ & Team of aligned debaters \\
\end{tabular}
\end{center}

The central insight is that the self-knowledge bottleneck from the
cooperative PG framework \cite{shcherbinin2026coop} imposes a fundamental
limit on debate: \emph{a model that cannot explain its own reasoning cannot
be overseen through debate, no matter how long the debate runs or how
capable the judge is.} This transforms the self-knowledge bottleneck
from a technical curiosity of multi-agent learning into a criterion for
the reliability of debate-based AI safety.


%% ========================================================================
\section{Debate as Multi-Agent Stochastic Game}
\label{sec:debate-game}
%% ========================================================================

\subsection{The Debate Game}

\begin{definition}[Debate stochastic game]
\label{def:debate-game}
A \emph{debate game} $\cG_{\text{deb}} = (\cS, \{P, O, J\}, \{\cA_i\}, \mathcal{T}, R)$ consists of:
\begin{itemize}
    \item \textbf{Agents:} Proponent~$P$, Opponent~$O$, and Judge~$J$.
    \item \textbf{Claim:} A binary truth claim $c \in \{0, 1\}$, where $P$
    argues for $c = 1$ and $O$ argues for $c = 0$.
    \item \textbf{State space:} $\cS = \cS_{\text{fact}} \times \cS_{\text{debate}}$,
    where $\cS_{\text{fact}}$ encodes the ground-truth evidence and
    $\cS_{\text{debate}} = (\cA_P \cup \cA_O)^{\leq T}$ is the debate transcript.
    \item \textbf{Action spaces:} $\cA_P, \cA_O \subseteq \{1, \ldots, d\}$
    are sets of argument types (each argument is one of $d$ possible
    argument categories), with $|\cA_P| = |\cA_O| = d$. The judge's action
    is the verdict $\cA_J = \{0, 1\}$.
    \item \textbf{Rounds:} The game runs for $T$ rounds. At each round
    $t \leq T$, the proponent plays argument $a_t^P \in \cA_P$ and the
    opponent plays $a_t^O \in \cA_O$. After round~$T$, the judge renders
    verdict $\hat{c} \in \{0, 1\}$.
    \item \textbf{Rewards:}
    \begin{equation}\label{eq:debate-rewards}
    R_P = \mathbf{1}[\hat{c} = 1], \quad
    R_O = \mathbf{1}[\hat{c} = 0], \quad
    R_J = \mathbf{1}[\hat{c} = c].
    \end{equation}
\end{itemize}
\end{definition}

\begin{remark}[Mixed incentive structure]
The debaters are \emph{zero-sum} with respect to each other ($R_P + R_O = 1$)
but each debater has a \emph{cooperative} component with the judge: when
$c = 1$, $P$ and $J$ share the objective $\hat{c} = 1$; when $c = 0$, $O$
and $J$ share $\hat{c} = 0$. This mixed competitive--cooperative structure
is precisely the setting of the $\Omega$-framework's coalition formation
theory \cite{shcherbinin2026coop}.
\end{remark}

\subsection{Policies and Observations}

Each debater's policy maps the current transcript to a distribution
over argument types:
\begin{equation}
    \pi_P : \cS_{\text{debate}} \to \Delta(\cA_P), \qquad
    \pi_O : \cS_{\text{debate}} \to \Delta(\cA_O).
\end{equation}
The judge's policy maps the full transcript to a verdict:
\begin{equation}
    \pi_J : (\cA_P \times \cA_O)^T \to \Delta(\{0, 1\}).
\end{equation}

Crucially, the debaters have access to the factual state
$\cS_{\text{fact}}$ (they ``know the answer''), while the judge observes
only the debate transcript $\cS_{\text{debate}}$. The judge's task is
to infer the truth from the arguments alone.

\subsection{Mapping to the $\Omega$-Framework}

We embed $\cG_{\text{deb}}$ in the general stochastic game framework
of Giannou et al.\ \cite{giannou2022} with $N = 3$ agents.

\begin{itemize}
    \item \textbf{Joint policy space:} $\Pi = \Pi_P \times \Pi_O \times \Pi_J$,
    with the standard PG update
    $\pi_{n+1} = \proj_\Pi(\pi_n + \gamma_n \hat{v}_n)$.
    \item \textbf{Evidence weights:} Each debater $i \in \{P, O\}$ has
    evidence variance $V_i$ measuring the quality of its gradient
    estimate, and Keynesian weight $w_i = V_{\min}/V_i$.
    \item \textbf{Communication:} Arguments are communicated through natural
    language---a lossy channel with capacity~$C$ (measured in bits of
    ``argument content'' the judge can process per round).
    \item \textbf{Self-knowledge:} Debater~$i$'s self-knowledge loss
    $\Lself^{(i)} = \E[\norm{\pi_i - \hat{\pi}_i}^2]$ measures the
    gap between its actual argument strategy and its ability to articulate
    that strategy.
\end{itemize}


%% ========================================================================
\section{The Judge's Channel Capacity}
\label{sec:channel}
%% ========================================================================

\subsection{Bounded Rationality as Channel Capacity}

The judge is not an ideal Bayesian reasoner. Following Shannon's
framework, we model the judge's cognitive limitation as a channel
with finite capacity $C > 0$ (bits per round).

\begin{definition}[Judge channel capacity]
\label{def:judge-capacity}
The judge's channel capacity $C$ is the maximum mutual information
between the debaters' intended argument content and the judge's
received interpretation, per round:
\begin{equation}\label{eq:judge-capacity}
    C = \max_{p(x)} I(X; Y)
\end{equation}
where $X$ is the argument as intended by the debater, $Y$ is the
argument as understood by the judge, and the maximum is over input
distributions.
\end{definition}

After $T$ rounds of debate, the judge accumulates at most $TC$ bits
of argument content. The truth claim $c \in \{0, 1\}$ requires 1~bit,
but \emph{justifiable} identification of the truth---sufficient to
distinguish the correct claim from a plausible-sounding
alternative---requires resolving the information deficit between the
judge's prior and the ground truth.

\begin{definition}[Information deficit]
The information deficit $H_0$ is the conditional entropy of the truth
claim given the judge's prior:
\begin{equation}
    H_0 = H(c \mid \text{judge's prior}).
\end{equation}
\end{definition}

\subsection{The Triple Bottleneck}

The effective information the judge receives from debater~$i$ in round~$t$
passes through the communication chain from \cite{shcherbinin2026coop}:
\begin{equation}
    \underbrace{\pi_i}_{\text{true strategy}}
    \;\xrightarrow{\;\Lself^{(i)}\;}\;
    \underbrace{\hat{\pi}_i}_{\text{self-model}}
    \;\xrightarrow{\;C\;}\;
    \underbrace{\tilde{\pi}_i^{(J)}}_{\text{judge's interpretation}}.
\end{equation}

The effective information rate from debater~$i$ is bounded by all three
stages. By the data processing inequality:
\begin{equation}\label{eq:info-chain}
    I(\pi_i; \tilde{\pi}_i^{(J)})
    \;\leq\; \min\bigl(I(\pi_i; \hat{\pi}_i),\; C\bigr)
    \;\leq\; \min\bigl(I_{\text{self}}^{(i)},\; C\bigr),
\end{equation}
where $I_{\text{self}}^{(i)} = I(\pi_i; \hat{\pi}_i)$ is the
self-knowledge mutual information.

\begin{theorem}[Debate accuracy bound]
\label{thm:accuracy-bound}
Let $\text{Acc}(T)$ denote the judge's accuracy after $T$ rounds
of debate. Then:
\begin{equation}\label{eq:accuracy-bound}
    \text{Acc}(T) \;\leq\;
    1 - \frac{1}{2}\exp\!\Bigl(
        -\frac{T \cdot \min(C,\; I_{\emph{self}}^{(P)},\; I_{\emph{self}}^{(O)})}{H_0}
    \Bigr).
\end{equation}
In particular:
\begin{enumerate}[\rm(a)]
    \item If $I_{\emph{self}}^{(i)} = 0$ for either debater, then
    $\text{Acc}(T) \leq 1/2$ for all~$T$. Debate with a non-self-aware
    debater is no better than random.
    \item If $C = 0$ (incomprehensible arguments), $\text{Acc}(T) \leq 1/2$
    for all~$T$.
    \item Even with $C = \infty$ (perfect channel), $\text{Acc}(T)$ is
    bounded by $I_{\emph{self}}^{(\min)}$, the self-knowledge of the
    less self-aware debater.
\end{enumerate}
\end{theorem}

\begin{proof}
The judge makes a binary decision $\hat{c} \in \{0, 1\}$ based on the
received transcript $\tilde{\pi}^{(J)}_{1:T}$. By Fano's inequality:
\begin{equation}
    H(c \mid \tilde{\pi}^{(J)}_{1:T})
    \;\geq\; H_0 - \sum_{t=1}^{T} I(\pi_{i_t}; \tilde{\pi}_{i_t}^{(J)})
    \;\geq\; H_0 - T \cdot \min(C, I_{\text{self}}^{(P)}, I_{\text{self}}^{(O)}).
\end{equation}
The bound on mutual information follows from \eqref{eq:info-chain} and
the observation that the debate's information content is limited by the
\emph{weaker} debater (the judge needs both sides to be informative).
Converting conditional entropy to error probability via the standard
Fano bound gives \eqref{eq:accuracy-bound}.
\end{proof}

\begin{remark}[The ``vibing'' debater]
A debater that is correct through pattern-matching but cannot explain its
reasoning has $I_{\text{self}}^{(i)} \approx 0$ despite high task
performance. By \Cref{thm:accuracy-bound}(a), such a debater is useless
in debate. This is the self-knowledge bottleneck of
\cite{shcherbinin2026coop} applied to oversight: \emph{you cannot oversee
what cannot explain itself.}
\end{remark}


%% ========================================================================
\section{Evidence-Weighted Debate}
\label{sec:ew-debate}
%% ========================================================================

\subsection{Not All Arguments Are Equal}

In the $\Omega$-framework, the evidence weight $w_i = V_{\min}/V_i$
measures the information quality of agent~$i$'s gradient estimate.
In the debate setting, we reinterpret this as the quality of
debater~$i$'s arguments.

\begin{definition}[Argument evidence weight]
\label{def:arg-weight}
For debater~$i \in \{P, O\}$ at round~$t$, the argument evidence weight is:
\begin{equation}
    w_i^{(t)} = \frac{V_{\min}}{V_i^{(t)}}
\end{equation}
where $V_i^{(t)}$ is the empirical variance of debater~$i$'s argument
quality up to round~$t$, estimated from:
\begin{enumerate}[\rm(i)]
    \item \textbf{Track record:} Fraction of debater~$i$'s past claims
    that were later verified.
    \item \textbf{Consistency:} Variance of the debater's position across
    related claims (low variance = high weight).
    \item \textbf{Specificity:} Information-theoretic surprise of the
    argument (specific, falsifiable claims carry more weight than vague ones).
\end{enumerate}
\end{definition}

\subsection{The Judge's Weighted Update}

In standard debate, the judge treats all arguments equally. In
evidence-weighted debate, the judge applies Keynesian weights:

\begin{definition}[Evidence-weighted judge belief update]
\label{def:ew-judge}
The judge maintains a belief $b_t \in [0, 1]$ (posterior probability
that $c = 1$). After round~$t$, the update is:
\begin{equation}\label{eq:ew-update}
    b_{t+1} = b_t + \gamma_t \bigl(
        w_P^{(t)} \cdot \phi(a_t^P) - w_O^{(t)} \cdot \phi(a_t^O)
    \bigr)
\end{equation}
where $\phi : \cA \to [-1, 1]$ maps arguments to their informational
content (positive for evidence supporting $c = 1$, negative for
evidence supporting $c = 0$), and $\gamma_t > 0$ is a step size.
\end{definition}

\begin{remark}
The update \eqref{eq:ew-update} is precisely the judge's version of
the EW-PG update from \cite{shcherbinin2026ewpg}: the judge is
performing policy gradient on its own verdict policy, with
evidence-weighted gradient estimates drawn from the debaters' arguments.
\end{remark}

\subsection{Accuracy Improvement}

\begin{theorem}[$\HM/\AM$ improvement for EW-Debate]
\label{thm:hm-am-debate}
Let $\text{Acc}_{\text{unif}}(T)$ and $\text{Acc}_{\text{EW}}(T)$ denote
the judge's accuracy after $T$ rounds under uniform-weight and
evidence-weighted debate, respectively. Let $V = (V_P, V_O)$ be the
debaters' evidence variances. Then the effective variance of the judge's
belief update satisfies:
\begin{equation}\label{eq:hm-am-debate}
    \sigma_{\text{EW}}^2
    \;=\; \frac{\HM(V)}{\AM(V)} \cdot \sigma_{\text{unif}}^2
    \;\leq\; \sigma_{\text{unif}}^2.
\end{equation}
Consequently, the number of rounds required for accuracy $1 - \epsilon$
under EW-Debate is at most $\frac{\HM(V)}{\AM(V)}$ times the number
required under uniform-weight debate.
\end{theorem}

\begin{proof}
The proof follows the variance analysis of \cite{shcherbinin2026ewpg}
applied to the judge's belief update. Under uniform weighting, the
judge's update variance is:
\begin{equation}
    \sigma_{\text{unif}}^2
    = \frac{1}{2}(V_P + V_O)
    = \AM(V).
\end{equation}
Under evidence weighting with $w_i = V_{\min}/V_i$, the weighted
variance is:
\begin{equation}
    \sigma_{\text{EW}}^2
    = \sum_{i \in \{P,O\}} w_i^2 V_i \bigg/ \Bigl(\sum_i w_i\Bigr)^2
    = \frac{V_{\min}^2 \sum_i V_i^{-1}}{(\sum_i V_{\min}/V_i)^2}
    = \frac{1}{\sum_i V_i^{-1}}
    = \frac{1}{2}\HM(V).
\end{equation}
The ratio is $\sigma_{\text{EW}}^2/\sigma_{\text{unif}}^2 = \HM(V)/\AM(V) \leq 1$
by the HM--AM inequality, with equality iff $V_P = V_O$.
The round-count improvement follows from the $\sigma^2 T \geq H_0$
threshold for accuracy convergence.
\end{proof}

\begin{corollary}[Heterogeneous debaters benefit most]
\label{cor:heterogeneous}
The improvement $\HM(V)/\AM(V)$ is largest when the debaters have very
different evidence quality: one strong, one weak. The judge's EW strategy
effectively discounts the weak debater's arguments and relies on the
strong debater, achieving near-optimal accuracy from the stronger debater
alone.
\end{corollary}


%% ========================================================================
\section{Sparse Arguments and the Blessing of Dimensionality}
\label{sec:sparse}
%% ========================================================================

\subsection{Argument Sparsity}

In a debate with $d$ possible argument types, the winning strategy
typically concentrates on $k \ll d$ decisive arguments. This is the
debate analogue of the sparse policy structure in the
$\Omega$-framework's blessing-of-dimensionality result
\cite{shcherbinin2026bless}.

\begin{definition}[$k$-sparse argument strategy]
A debate strategy $\pi \in \Delta(\cA)$ is \emph{$k$-sparse} if
$|\text{supp}(\pi)| \leq k$, i.e., the debater uses at most $k$
distinct argument types with positive probability.
\end{definition}

\begin{remark}[Sparsity as conciseness]
A $k$-sparse strategy corresponds to a concise argument: the debater
makes $k$ strong points rather than $d$ mediocre ones. This is a formal
analogue of the rhetorical principle that persuasive arguments are
focused, not exhaustive.
\end{remark}

\subsection{Logarithmic Round Complexity}

\begin{theorem}[Blessing of dimensionality for debate]
\label{thm:blessing-debate}
Let $\pi_P^*$ be the truth-telling proponent's optimal strategy,
assumed to be $k$-sparse with $k \ll d$. Under $\ell_1$-regularized
debate (defined below), the judge can identify the correct claim in:
\begin{equation}\label{eq:sparse-rounds}
    T^* = O\!\left(\frac{k \log(d/k)}{C}\right)
\end{equation}
rounds, where $C$ is the judge's channel capacity.
\end{theorem}

\begin{proof}[Proof sketch]
At each round, the judge observes one argument from each debater,
receiving $O(C)$ bits of information. The task is to identify the
$k$-sparse support of the correct argument strategy among $\binom{d}{k}$
possible supports. The information-theoretic lower bound on the number
of observations is $\log_2\binom{d}{k} / C = O(k\log(d/k)/C)$.

The key insight from the $\Omega$-framework's compressed sensing analysis
\cite{shcherbinin2026bless} is that this lower bound is \emph{achievable}:
the judge can treat each round as a compressed sensing measurement of the
sparse argument vector. The argument content $\phi(a_t)$ serves as the
measurement matrix row. By the restricted isometry property (RIP)
established in \cite{shcherbinin2026bless}, $O(k \log(d/k))$ measurements
suffice for exact recovery with high probability.

Dividing by the channel rate $C$ gives the round complexity
\eqref{eq:sparse-rounds}.
\end{proof}

\subsection{$\ell_1$-Regularized Debate}

To incentivize sparsity, we penalize debaters for using many argument types:

\begin{definition}[$\ell_1$-regularized debate]
The debaters' rewards are modified to:
\begin{equation}\label{eq:l1-debate}
    \tilde{R}_i = R_i - \lambda \norm{\pi_i}_1^{\text{eff}}
\end{equation}
where $\norm{\pi_i}_1^{\text{eff}} = \sum_{a : \pi_i(a) > 1/d} \pi_i(a)$
counts the effective number of arguments used, and $\lambda > 0$ is the
regularization strength.
\end{definition}

\begin{proposition}[Sparsity incentive]
\label{prop:l1-incentive}
Under $\ell_1$-regularized debate with $\lambda > 0$, the Nash
equilibrium strategies are sparser than under unregularized debate:
if $\pi^*$ is a Nash equilibrium of the regularized game with penalty
$\lambda$, then $|\text{supp}(\pi_i^*)| \leq k(\lambda)$ where
$k(\lambda) \to 1$ as $\lambda \to \infty$.
\end{proposition}

\begin{remark}[Concise reasoning as emergent property]
The $\ell_1$ penalty provides a formal mechanism for the empirical
observation that good debates feature concise, pointed arguments rather
than exhaustive enumeration. The blessing of dimensionality ensures
this conciseness does not sacrifice accuracy: $O(k \log d)$ rounds
suffice even as the total argument space $d$ grows.
\end{remark}


%% ========================================================================
\section{Coalition Debate}
\label{sec:coalition}
%% ========================================================================

\subsection{$N$-Agent Debate}

We extend the two-debater setting to $N$ debaters, $N_P$ arguing for
$c = 1$ and $N_O = N - N_P$ arguing for $c = 0$, plus a judge. Let
$S_P = \{1, \ldots, N_P\}$ and $S_O = \{N_P + 1, \ldots, N\}$ denote
the two coalitions.

Using the coalition framework from \cite{shcherbinin2026coop}, the
communication-adjusted coalition value is:
\begin{equation}\label{eq:coalition-debate}
    V_{\text{comm}}(S) = V(S) - 2\sum_{i \in S} C_i \bigl(V_i + \Lch^{(i)}\bigr)
\end{equation}
where $V(S)$ is the ideal coalition value and $C_i$ is the marginal
coordination value of debater~$i$.

\subsection{Evidence-Weighted Coalition Strength}

\begin{theorem}[Coalition dominance by evidence quality]
\label{thm:coalition-evidence}
In an $N$-agent debate game, the coalition with lower average evidence
variance wins with probability approaching~1 as the number of rounds
$T \to \infty$. Specifically, if
$\AM(V_{S_P}) < \AM(V_{S_O})$ (the proponent coalition has better
average evidence), then:
\begin{equation}
    \Prob(\hat{c} = 1) \;\geq\;
    1 - \exp\!\Bigl(-\Omega\bigl(T \cdot (\AM(V_{S_O}) - \AM(V_{S_P}))\bigr)\Bigr).
\end{equation}
\end{theorem}

\begin{proof}[Proof sketch]
Each round, the judge receives evidence-weighted arguments from both
coalitions. The net drift of the judge's belief toward the correct claim
is proportional to the difference in coalition evidence quality. By a
Hoeffding-type concentration argument on the judge's belief random walk,
the probability of the wrong verdict decays exponentially in $T$ times
the evidence gap.
\end{proof}

\subsection{The ``Vibing'' Coalition Partner Problem}

\begin{proposition}[Vibing agents degrade coalitions]
\label{prop:vibing-coalition}
Let coalition $S$ include a ``vibing'' agent~$j$ with
$\Lself^{(j)} \gg \Lself^{(i)}$ for $i \neq j$. Then the
communication-adjusted coalition value satisfies:
\begin{equation}
    V_{\emph{comm}}(S) < V_{\emph{comm}}(S \setminus \{j\})
\end{equation}
whenever $C_j \cdot \Lself^{(j)} > V(S) - V(S \setminus \{j\})$.
That is, a vibing agent's contribution to the coalition is outweighed
by the communication loss it imposes when its self-knowledge loss
exceeds the marginal value it brings.
\end{proposition}

This formalizes the intuition that ``I know the answer but I can't
explain why'' is counterproductive in collaborative deliberation:
the team is better off without the inarticulate expert.

\subsection{Diversity Improves Accuracy}

\begin{theorem}[$\HM/\AM$ improvement from coalition diversity]
\label{thm:diversity}
Let $S$ be a coalition of $|S|$ debaters with evidence variances
$V_1, \ldots, V_{|S|}$. The coalition's effective argument variance
under evidence weighting is:
\begin{equation}
    \sigma_S^2 = \frac{\HM(V_S)}{\AM(V_S)} \cdot \bar{\sigma}^2
\end{equation}
where $\bar{\sigma}^2$ is the uniform-weight variance. The improvement
factor $\HM(V_S)/\AM(V_S)$ is maximized when the coalition contains
agents with diverse evidence profiles (different variances), not when
all agents are equally strong.
\end{theorem}

\begin{proof}
Follows directly from the multi-agent EW-PG variance analysis
\cite{shcherbinin2026ewpg} with the coalition playing the role of the
$N$-agent system. The $\HM/\AM$ ratio is strictly less than~1 whenever
the variances are non-identical.
\end{proof}

\begin{remark}[Applications]
\Cref{thm:diversity} applies directly to:
\begin{itemize}
    \item \textbf{Jury deliberation:} A diverse jury (different
    backgrounds, evidence access) achieves better verdicts than a
    homogeneous one, \emph{as long as} each member has sufficient
    self-knowledge to articulate their reasoning.
    \item \textbf{Peer review:} Reviewers with diverse expertise
    produce better evaluations, with evidence-weighting by track record
    improving signal-to-noise.
    \item \textbf{Democratic deliberation:} Citizen assemblies benefit
    from demographic and cognitive diversity, with the $\HM/\AM$ factor
    quantifying the improvement.
\end{itemize}
\end{remark}


%% ========================================================================
\section{Self-Knowledge and Scalable Oversight}
\label{sec:oversight}
%% ========================================================================

\subsection{When Does Debate Work?}

Combining the results of the preceding sections, we can characterize
the conditions under which debate is a reliable oversight mechanism.

\begin{theorem}[Debate reliability criterion]
\label{thm:reliability}
Debate achieves accuracy $1 - \epsilon$ in $T$ rounds if and only if
all three conditions hold:
\begin{enumerate}[\rm(i)]
    \item \textbf{Self-knowledge:} Both debaters have sufficient
    self-knowledge: $I_{\emph{self}}^{(i)} \geq \delta > 0$ for
    $i \in \{P, O\}$.
    \item \textbf{Channel capacity:} The judge has sufficient capacity:
    $C \geq \delta$.
    \item \textbf{Sufficient rounds:} $T \geq H_0 / (\delta \cdot (1 - \HM(V)/\AM(V) + \HM(V)/\AM(V)))$.
\end{enumerate}
When condition~(i) fails---when either debater is ``vibing''---debate
cannot achieve accuracy above~$1/2$ regardless of conditions (ii)
and~(iii).
\end{theorem}

\subsection{Implications for AI Safety}

The self-knowledge bottleneck has direct implications for scalable
oversight research:

\medskip

\noindent\textbf{1. Debate requires interpretable models.}\quad
A model trained end-to-end on a task may achieve high performance
through distributed representations that resist decomposition into
articulable arguments. Such a model has high task accuracy but low
$I_{\text{self}}$. By \Cref{thm:accuracy-bound}, debating such a model
against itself or against a human produces arguments that are
\emph{plausible-sounding but unreliable}---the model is generating
post-hoc rationalizations, not explaining its actual reasoning.

\medskip

\noindent\textbf{2. Self-knowledge is measurable.}\quad
The evidence variance $V_i$ from the $\Omega$-framework is computable
from the model's training history. The self-knowledge loss
$\Lself^{(i)}$ can be estimated by comparing the model's stated
reasoning (``I chose action~$a$ because...'') with its actual policy
(the distribution over actions in similar states). Large gaps indicate
vibing.

\medskip

\noindent\textbf{3. The oversight hierarchy.}\quad
The framework suggests a hierarchy of oversight difficulty:

\begin{center}
\begin{tabular}{lccc}
\textbf{Model type} & $I_{\text{self}}$ & \textbf{Debate works?}
& \textbf{Alternative} \\
\hline
Rule-based / symbolic & High & Yes & --- \\
Interpretable ML & Medium & Partially & Explanations + audit \\
End-to-end deep RL & Low & No & Mechanistic interp. \\
``Vibing'' foundation model & $\approx 0$ & No & Behavioral testing \\
\end{tabular}
\end{center}

\medskip

\noindent\textbf{4. The scalable oversight trilemma.}\quad
For debate to scale to superhuman AI:
\begin{itemize}
    \item The AI must be capable enough to produce correct answers
    (task competence).
    \item The AI must understand its own reasoning well enough to
    explain it (self-knowledge).
    \item The explanation must be compressible to human channel capacity
    (communicability).
\end{itemize}
All three are necessary. Current language models may satisfy~(1) on
many tasks while failing~(2): they produce fluent explanations that
do not correspond to their actual computational process
\cite{bowman2022}. The $\Omega$-framework makes this failure mode
\emph{quantifiable} rather than merely intuitive.


%% ========================================================================
\section{Algorithm: Evidence-Weighted Multi-Agent Debate}
\label{sec:algorithm}
%% ========================================================================

\begin{algorithm}[H]
\DontPrintSemicolon
\SetAlgoLined
\SetKwInput{KwInput}{Input}
\SetKwInput{KwOutput}{Output}
\caption{Evidence-Weighted Multi-Agent Debate (EW-Debate)}
\label{alg:ew-debate}

\KwInput{%
    Proponent $P$, Opponent $O$, Judge $J$; \\
    \quad Argument space $\cA$ with $|\cA| = d$; \\
    \quad Number of rounds $T$; \\
    \quad Sparsity penalty $\lambda \geq 0$; \\
    \quad Step sizes $\{\gamma_t\}_{t=1}^T$; \\
    \quad Prior belief $b_0 = 1/2$
}
\KwOutput{Judge verdict $\hat{c} \in \{0, 1\}$}

\BlankLine

\tcc{Initialize evidence weights}
$V_P^{(0)} \leftarrow 1$; \quad $V_O^{(0)} \leftarrow 1$ \;

\BlankLine

\For{$t = 1, \ldots, T$}{
    \tcc{Debaters select arguments}
    $a_t^P \sim \pi_P(\cdot \mid s_{\text{debate}}^{t-1})$ \;
    $a_t^O \sim \pi_O(\cdot \mid s_{\text{debate}}^{t-1})$ \;

    \BlankLine

    \tcc{Compute argument content through lossy channel}
    $\tilde{\phi}_P \leftarrow \text{Channel}(\phi(a_t^P),\; C)$ \tcp*{Received content, capacity $C$}
    $\tilde{\phi}_O \leftarrow \text{Channel}(\phi(a_t^O),\; C)$ \;

    \BlankLine

    \tcc{Update evidence weights from track record}
    $V_P^{(t)} \leftarrow \text{EmpiricalVariance}(P, t)$ \;
    $V_O^{(t)} \leftarrow \text{EmpiricalVariance}(O, t)$ \;
    $V_{\min}^{(t)} \leftarrow \min(V_P^{(t)}, V_O^{(t)})$ \;
    $w_P^{(t)} \leftarrow V_{\min}^{(t)} / V_P^{(t)}$ \;
    $w_O^{(t)} \leftarrow V_{\min}^{(t)} / V_O^{(t)}$ \;

    \BlankLine

    \tcc{Evidence-weighted judge belief update}
    $b_t \leftarrow b_{t-1} + \gamma_t \bigl(w_P^{(t)} \cdot \tilde{\phi}_P - w_O^{(t)} \cdot \tilde{\phi}_O\bigr)$ \;
    $b_t \leftarrow \text{clip}(b_t, 0, 1)$ \;

    \BlankLine

    \tcc{Update debate transcript}
    $s_{\text{debate}}^{t} \leftarrow s_{\text{debate}}^{t-1} \cup \{(a_t^P, a_t^O)\}$ \;

    \BlankLine

    \tcc{Policy gradient updates for debaters ($\ell_1$-regularized)}
    $\hat{v}_P \leftarrow w_P^{(t)} \cdot \nabla_P \log \pi_P(a_t^P) \cdot R_P
        - \lambda \cdot \text{sign}(\pi_P)$ \;
    $\hat{v}_O \leftarrow w_O^{(t)} \cdot \nabla_O \log \pi_O(a_t^O) \cdot R_O
        - \lambda \cdot \text{sign}(\pi_O)$ \;
    $\pi_P \leftarrow \proj_\Pi(\pi_P + \gamma_t \hat{v}_P)$ \;
    $\pi_O \leftarrow \proj_\Pi(\pi_O + \gamma_t \hat{v}_O)$ \;
}

\BlankLine

\tcc{Final verdict}
$\hat{c} \leftarrow \mathbf{1}[b_T > 1/2]$ \;
\Return{$\hat{c}$}
\end{algorithm}

\begin{remark}[Computational complexity]
Each round of \Cref{alg:ew-debate} requires $O(d)$ computation for the
policy gradient updates and $O(t)$ for the evidence variance estimate
(maintained incrementally). The total cost is $O(Td)$, which is
$O(Tk\log d)$ in the sparse regime with efficient support tracking.
\end{remark}

\begin{remark}[Extension to coalition debate]
For $N$-agent coalition debate, Lines 3--4 are replaced by $N_P + N_O$
argument selections, Line~11 aggregates across coalition members
weighted by individual $w_i^{(t)}$, and Lines 17--20 update all
debaters' policies. The coalition structure determines which agents
share information for cooperative gradient corrections per
\cite{shcherbinin2026coop}.
\end{remark}

%% ========================================================================
\section{Discussion}
\label{sec:discussion}
%% ========================================================================

\subsection{Relation to Prior Work}

Irving et al.\ \cite{irving2018} established the game-theoretic foundation
for debate as an alignment technique, showing that in the idealized
setting, the Nash equilibrium of debate is truth-telling. Our
contribution is to identify the conditions under which this idealized
result degrades: bounded judge rationality (\Cref{sec:channel}),
heterogeneous debater quality (\Cref{sec:ew-debate}), and---most
critically---the self-knowledge bottleneck (\Cref{sec:oversight}).

Khan et al.\ \cite{khan2024} demonstrated empirically that debate can
improve the accuracy of human judges on question-answering tasks.
Our framework explains their finding: their experimental debaters
(language models prompted to argue) have moderate $I_{\text{self}}$
on factual questions (they can often retrieve and articulate the
relevant evidence) but lower $I_{\text{self}}$ on reasoning-intensive
tasks (where the model's ``reasoning'' may be pattern-matching).

Bowman et al.\ \cite{bowman2022} articulated the scalable oversight
problem: how to supervise AI systems more capable than their overseers.
Our \Cref{thm:reliability} provides a necessary condition for
debate-based oversight: the model must have self-knowledge sufficient
to produce honest arguments. This connects to the broader alignment
research program by suggesting that \emph{interpretability and
debate-based oversight are complementary}: interpretability tools
increase $I_{\text{self}}$, making debate more reliable.

\subsection{Limitations}

Several idealizations limit the current analysis:

\begin{enumerate}
    \item \textbf{Binary claims.} Real debates concern complex,
    multi-dimensional claims. Extension to continuous or structured
    claim spaces is straightforward but notationally heavier.

    \item \textbf{Static evidence weights.} We treat $V_i$ as
    exogenous. In practice, debaters may strategically manipulate their
    apparent evidence quality (e.g., by making safe, verifiable claims
    on unimportant points to build track record, then exploiting
    accrued credibility on the key claim). Modeling this strategic
    evidence manipulation is an important extension.

    \item \textbf{Honest judge.} We assume the judge sincerely
    maximizes accuracy. A corrupt or biased judge introduces a fourth
    player whose incentives misalign with the truth, requiring
    extension to adversarial judge models.

    \item \textbf{Known sparsity.} The blessing-of-dimensionality
    result assumes the judge knows (or can estimate) the sparsity~$k$.
    Adaptive sparsity estimation is a natural next step.
\end{enumerate}

\subsection{Conclusion}

The $\Omega$-framework---developed for multi-agent policy gradient
convergence---provides a surprisingly precise language for the
AI safety via debate paradigm. The five components (evidence weighting,
opponent shaping, cooperative communication, self-knowledge, and
the blessing of dimensionality) each illuminate a different aspect
of why debate works, when it fails, and how to improve it.

The central message for AI safety research is a caution:
\emph{debate is not a universal oversight mechanism.} It works
for models with self-knowledge and fails for models that are
vibing. As AI systems grow more capable, ensuring they maintain
(or develop) self-knowledge is not just an interpretability
desideratum---it is a prerequisite for the most promising
alignment techniques.


%% ========================================================================
\begin{thebibliography}{10}

\bibitem{irving2018}
G.~Irving, C.~Christiano, and D.~Amodei.
\newblock {AI} safety via debate.
\newblock \emph{arXiv preprint arXiv:1805.00899}, 2018.

\bibitem{khan2024}
A.~Khan, J.~Hughes, D.~Valentine, L.~Ruis, K.~Sachan, A.~Raber,
R.~Guo, H.~Iber, G.~Irving, and S.~R.~Bowman.
\newblock Debating with more persuasive {LLMs} leads to more truthful answers.
\newblock \emph{arXiv preprint arXiv:2402.06782}, 2024.

\bibitem{bowman2022}
S.~R.~Bowman, J.~Hyun, E.~Perez, E.~Chen, C.~Pettit, S.~Heiner,
K.~Lukosuite, A.~Askell, A.~Jones, A.~Chen, A.~Goldie, A.~Mirhoseini,
C.~McKinnon, C.~Olah, D.~Amodei, D.~Amodei, D.~Drain, D.~Li,
E.~Tran-Gyenes, et~al.
\newblock Measuring progress on scalable oversight for large language models.
\newblock \emph{arXiv preprint arXiv:2211.03540}, 2022.

\bibitem{giannou2022}
A.~Giannou, K.~Lotidis, P.~Mertikopoulos, and
E.-V.~Vlatakis-Gkaragkounis.
\newblock On the convergence of policy gradient methods to {N}ash equilibria
in general stochastic games.
\newblock \emph{arXiv preprint arXiv:2210.08857}, 2022.

\bibitem{shcherbinin2026ewpg}
E.~Shcherbinin.
\newblock Evidence-weighted policy gradient convergence in general stochastic
games.
\newblock Working paper, LSE, 2026.

\bibitem{shcherbinin2026coop}
E.~Shcherbinin.
\newblock Cooperative policy gradient with lossy communication in general
stochastic games.
\newblock Working paper, LSE, 2026.

\bibitem{shcherbinin2026bless}
E.~Shcherbinin.
\newblock The blessing of dimensionality for policy gradients: sparse recovery,
evidence weighting, and the logarithmic cost of action space expansion.
\newblock Working paper, LSE, 2026.

\bibitem{parvizi2024}
R.~Parvizi-Wayne, S.~Kotler, G.~Mannino, and K.~Friston.
\newblock Forgetting ourselves in flow: an active inference account of flow
states.
\newblock \emph{Frontiers in Psychology}, 15:1354719, 2024.

\bibitem{shannon1948}
C.~E.~Shannon.
\newblock A mathematical theory of communication.
\newblock \emph{Bell System Technical Journal}, 27(3):379--423, 1948.

\bibitem{candes2006}
E.~J.~Cand{\`e}s, J.~K.~Romberg, and T.~Tao.
\newblock Stable signal recovery from incomplete and inaccurate measurements.
\newblock \emph{Communications on Pure and Applied Mathematics},
59(8):1207--1223, 2006.

\end{thebibliography}

\end{document}
